\section*{Chi tiết các Thực thể Dữ liệu}

% =========================================================
% A. NHÓM QUẢN TRỊ & KẾT NỐI
% =========================================================

\section{Nhóm Quản trị \& Kết nối}

\subsection{User}

\textbf{User (Thực thể mạnh):} Người dùng cuối trong hệ thống
(Developer, Staff, Manager, Tech Lead, Admin).

\subsubsection*{Thuộc tính Phức hợp (Composite Attribute)}

\begin{itemize}
    \item \textbf{FullName}: Phức hợp, gồm \textbf{FirstName} và \textbf{LastName}.
\end{itemize}

\subsubsection*{Thuộc tính Đa trị (Multi-valued Attribute)}

\begin{itemize}
    \item \textbf{NotificationPreferences}: Cài đặt thông báo, lưu nhiều giá trị cùng lúc cho một người dùng (ví dụ: chọn đồng thời cả EMAIL và SYSTEM).
\end{itemize}

\subsubsection*{Thuộc tính Đơn (Simple / Atomic Attribute)}

\begin{itemize}
    \item \textbf{UserID (PK)}: Định danh duy nhất cho từng người dùng.
    \item \textbf{Email}: Địa chỉ email bắt buộc dùng để đăng nhập và nhận thông báo hệ thống.
    \item \textbf{PasswordHash}: Mật khẩu đã mã hóa. Lưu giá trị khi đăng nhập bằng Email/Password và giữ \textbf{NULL} khi dùng thuần GitHub OAuth.
    \item \textbf{Status}: Trạng thái tài khoản (ACTIVE, INACTIVE, BLOCKED).
    \item \textbf{CreatedAt}: Thời điểm khởi tạo tài khoản.
    \item \textbf{UpdatedAt}: Thời điểm cập nhật thông tin gần nhất.
    \item \textbf{DeletedAt}: Soft delete nhằm giữ nguyên lịch sử giao dịch (review, comment, approval) khi người dùng rời hệ thống.
\end{itemize}


% =========================================================
% WORKSPACE
% =========================================================

\subsection{Workspace}

\textbf{Workspace (Thực thể mạnh):} Không gian làm việc chứa repository và tài liệu của một nhóm/dự án.

\subsubsection*{Thuộc tính Đơn (Simple / Atomic Attribute)}

\begin{itemize}
    \item \textbf{workspace\_id (PK)}: Định danh duy nhất cho từng không gian làm việc.
    \item \textbf{name}: Tên không gian làm việc/dự án (VD: ``Payment Service Platform'').
    \item \textbf{description}: Mô tả chi tiết về phạm vi và mục tiêu của workspace.
    \item \textbf{status}: Trạng thái hoạt động của workspace (ACTIVE, ARCHIVED).
    \item \textbf{created\_at}: Thời điểm khởi tạo workspace.
    \item \textbf{updated\_at}: Thời điểm cập nhật thông tin workspace gần nhất.
    \item \textbf{deleted\_at}: Cờ xóa mềm (Soft delete) giữ nguyên dữ liệu báo cáo và audit trail khi ẩn workspace.
\end{itemize}


% =========================================================
% ROLE
% =========================================================

\subsection{Role}

\textbf{Role (Thực thể mạnh):} Vai trò của người dùng (scoped theo Workspace).

\subsubsection*{Thuộc tính Đơn (Simple / Atomic Attribute)}

\begin{itemize}
    \item \textbf{role\_id (PK)}: Định danh duy nhất cho vai trò.
    \item \textbf{role\_name}: Tên vai trò (VD: Developer, Staff, Manager, Tech Lead, Admin).
    \item \textbf{description}: Mô tả tóm tắt mục đích và quyền hạn chung của vai trò.
    \item \textbf{created\_at}: Thời điểm khởi tạo vai trò trong hệ thống.
    \item \textbf{updated\_at}: Thời điểm cập nhật vai trò gần nhất.
\end{itemize}


% =========================================================
% PERMISSION
% =========================================================

\subsection{Permission}

\textbf{Permission (Thực thể mạnh):} Quyền hạn chi tiết trên hệ thống
(VD: DOC\_APPROVE, TEMPLATE\_EDIT).

\subsubsection*{Thuộc tính Đơn (Atomic Attribute)}

\begin{itemize}
    \item \textbf{permission\_id (PK)}: Định danh duy nhất cho quyền.
    \item \textbf{permission\_code}: Mã quyền thao tác hệ thống (VD: DOC\_APPROVE, TEMPLATE\_EDIT).
    \item \textbf{description}: Giải thích chi tiết giới hạn của thao tác này.
    \item \textbf{created\_at}: Mốc thời gian quyền được khởi tạo vào danh mục hệ thống.
\end{itemize}


\textbf{RolePermission (Thực thể liên kết):} Bản ghi phân quyền chi tiết cho từng vai trò trong hệ thống.

\textbf{-- Thuộc tính Đơn (Atomic Attribute):}

\begin{itemize}
    \item \textbf{role\_permission\_id (PK):} Định danh duy nhất cho bản ghi phân quyền (UUID).
    
    \item \textbf{role\_id (FK):} Khóa ngoại trỏ đến vai trò được gán quyền (\texttt{roles}).
    
    \item \textbf{permission\_id (FK):} Khóa ngoại trỏ đến quyền hạn cụ thể (\texttt{permissions}).
    
    \item \textbf{created\_at:} Mốc thời gian quyền này được gán cho vai trò.
\end{itemize}


% =========================================================
% WORKSPACE MEMBER
% =========================================================

\subsection{WorkspaceMember}

\textbf{WorkspaceMember (Thực thể mối kết hợp):} Ghi nhận sự tham gia của User trong Workspace và vai trò (Role) tương ứng.

\subsubsection*{Thuộc tính Đơn (Atomic Attribute)}

\begin{itemize}
    \item \textbf{member\_id (PK)}: Định danh duy nhất cho bản ghi tham gia.
    \item \textbf{status}: Trạng thái thành viên (ACTIVE, INACTIVE, REMOVED).
    \item \textbf{joined\_at}: Thời điểm người dùng chính thức tham gia.
    \item \textbf{created\_at}: Thời điểm gửi lời mời / tạo bản ghi.
    \item \textbf{updated\_at}: Thời điểm cập nhật vai trò hoặc trạng thái gần nhất.
\end{itemize}


% =========================================================
% APPROVAL DELEGATION
% =========================================================

\subsection{ApprovalDelegation}

\textbf{ApprovalDelegation (Thực thể phụ thuộc):} Bản ghi ủy quyền phê duyệt giữa các Manager trong khoảng thời gian xác định.

\subsubsection*{Thuộc tính Đơn (Simple / Atomic Attribute)}

\begin{itemize}
    \item \textbf{delegation\_id (PK)}: Định danh duy nhất cho bản ghi ủy quyền.
    \item \textbf{valid\_from}: Thời điểm bắt đầu có hiệu lực ủy quyền.
    \item \textbf{valid\_until}: Thời điểm kết thúc hiệu lực ủy quyền.
    \item \textbf{reason}: Lý do ủy quyền (VD: Nghỉ phép, công tác).
    \item \textbf{status}: Trạng thái ủy quyền (ACTIVE, EXPIRED, REVOKED).
    \item \textbf{created\_at}: Thời điểm tạo bản ghi ủy quyền.
    \item \textbf{updated\_at}: Thời điểm cập nhật trạng thái ủy quyền gần nhất.
\end{itemize}


% =========================================================
% GITHUB CONNECTION
% =========================================================

\subsection{GitHubConnection}

\textbf{GitHubConnection (Thực thể phụ thuộc):} Lưu thông tin xác thực OAuth của User với GitHub.

\subsubsection*{Thuộc tính Khóa (Key Attribute)}

\begin{itemize}
    \item \textbf{connection\_id}: Định danh duy nhất cho bản ghi kết nối OAuth.
\end{itemize}

\subsubsection*{Thuộc tính Đơn / Nguyên tố (Simple / Atomic Attributes)}

\begin{itemize}
    \item \textbf{github\_user\_id}: ID tài khoản người dùng trên hệ thống GitHub (để đối soát định danh).
    \item \textbf{github\_username}: Tên tài khoản GitHub (VD: octocat).
    \item \textbf{access\_token}: Chuỗi token ủy quyền dùng để gọi GitHub API.
    \item \textbf{created\_at}: Thời điểm người dùng thực hiện liên kết tài khoản.
    \item \textbf{updated\_at}: Thời điểm làm mới token hoặc cập nhật quyền hạn gần nhất.
\end{itemize}

\subsubsection*{Thuộc tính Đa trị (Multi-valued Attribute)}

\begin{itemize}
    \item \textbf{scopes}: Lưu danh sách các phạm vi quyền OAuth được cấp phép cùng lúc (VD: repo, read:user, workflow).
\end{itemize}


% =========================================================
% B. NHÓM MÃ NGUỒN & CẤU TRÚC CODE
% =========================================================

\section{Nhóm Mã nguồn \& Cấu trúc Code}

\subsection{Repository}

\textbf{Repository (Thực thể mạnh):} Kho chứa mã nguồn được đồng bộ vào hệ thống.

\subsubsection*{Thuộc tính Đơn (Atomic Attribute)}

\begin{itemize}
    \item \textbf{repository\_id (PK)}: Định danh duy nhất cho kho mã nguồn.
    \item \textbf{name}: Tên kho mã nguồn (VD: payment-gateway).
    \item \textbf{url}: Đường dẫn gốc đến kho chứa trên VCS.
    \item \textbf{default\_branch}: Nhánh mặc định dùng để phân tích code (VD: main, master).
    \item \textbf{sync\_status}: Trạng thái đồng bộ dữ liệu (SYNCING, SUCCESS, FAILED).
    \item \textbf{last\_synced\_at}: Thời điểm đồng bộ mã nguồn thành công gần nhất.
    \item \textbf{created\_at}: Thời điểm tích hợp kho mã nguồn vào hệ thống.
    \item \textbf{updated\_at}: Thời điểm cập nhật trạng thái hoặc cấu hình kho gần nhất.
\end{itemize}


% =========================================================
% BRANCH
% =========================================================

\subsection{Branch}

\textbf{Branch (Thực thể yếu -- phụ thuộc Repository):} Nhánh phát triển trong repository.

\subsubsection*{Thuộc tính Đơn (Simple / Atomic Attribute)}

\begin{itemize}
    \item \textbf{branch\_id (PK)}: Định danh duy nhất cho nhánh.
    \item \textbf{branch\_name}: Tên nhánh (VD: main, staging, feature/auth).
    \item \textbf{is\_default}: Cờ xác định nhánh mặc định của repo (true/false).
    \item \textbf{created\_at}: Thời điểm nhánh được nhận diện và lưu vào hệ thống.
    \item \textbf{updated\_at}: Thời điểm cập nhật commit hash hoặc trạng thái nhánh gần nhất.
\end{itemize}


% =========================================================
% COMMIT
% =========================================================

\subsection{Commit}

\textbf{Commit (Thực thể yếu -- phụ thuộc Branch/Repo):} Bản ghi một lần commit mã nguồn.

\subsubsection*{Thuộc tính Phức hợp (Composite Attribute)}

\begin{itemize}
    \item \textbf{AuthorInfo}: Thông tin tác giả thực hiện commit, gồm \textbf{AuthorName} và \textbf{AuthorEmail}.
\end{itemize}

\subsubsection*{Thuộc tính Đơn (Simple / Atomic Attribute)}

\begin{itemize}
    \item \textbf{commit\_id (PK)}: Định danh nội bộ cho bản ghi commit.
    \item \textbf{commit\_hash}: Mã băm SHA-1/SHA-256 duy nhất từ Git (VD: 7a8f9b2...).
    \item \textbf{message}: Nội dung mô tả của commit.
    \item \textbf{committed\_at}: Thời điểm commit được thực hiện trên VCS (thời gian nghiệp vụ).
    \item \textbf{created\_at}: Thời điểm bản ghi commit được đồng bộ vào LivingDocs (thời gian hệ thống).
\end{itemize}


% =========================================================
% PULL REQUEST
% =========================================================

\subsection{PullRequest}

\textbf{PullRequest (Thực thể yếu -- phụ thuộc Repository):} Yêu cầu hợp nhất nhánh (PR) từ GitHub.

\subsubsection*{Thuộc tính Đa trị (Multi-valued Attribute)}

\begin{itemize}
    \item \textbf{labels}: Danh sách nhãn phân loại được gắn trên PR (VD: bug, enhancement, documentation, api-change).
\end{itemize}

\subsubsection*{Thuộc tính Đơn (Simple / Atomic Attribute)}

\begin{itemize}
    \item \textbf{pr\_id (PK)}: Định danh nội bộ của bản ghi PR.
    \item \textbf{pr\_number}: Số thứ tự PR trên nền tảng gốc (VD: \#102, \#123).
    \item \textbf{title}: Tiêu đề của Pull Request.
    \item \textbf{description}: Nội dung văn bản mô tả chi tiết PR (Dữ liệu đầu vào cho AI).
    \item \textbf{source\_branch}: Nhánh nguồn chứa mã nguồn thay đổi.
    \item \textbf{target\_branch}: Nhánh đích dự định hợp nhất vào (VD: main, develop).
    \item \textbf{status}: Trạng thái vòng đời của PR (OPEN, MERGED, CLOSED).
    \item \textbf{created\_at}: Thời điểm mở PR trên VCS.
    \item \textbf{updated\_at}: Thời điểm PR cập nhật thông tin/trạng thái gần nhất.
    \item \textbf{merged\_at}: Thời điểm PR chính thức được chấp nhận hợp nhất (Nullable).
\end{itemize}


% =========================================================
% CODE DIFF
% =========================================================

\subsection{CodeDiff}

\textbf{CodeDiff (Thực thể yếu -- phụ thuộc Commit/PR):} Chi tiết thay đổi mã nguồn trong commit hoặc PR.

\subsubsection*{Thuộc tính Đơn (Simple / Atomic Attribute)}

\begin{itemize}
    \item \textbf{diff\_id (PK)}: Định danh nội bộ bản ghi chi tiết thay đổi.
    \item \textbf{file\_path}: Đường dẫn tương đối đến tệp bị thay đổi (VD: src/controllers/payment.js).
    \item \textbf{change\_type}: Loại thao tác trên file (ADD, MODIFY, DELETE, RENAME).
    \item \textbf{additions}: Số dòng code được thêm mới.
    \item \textbf{deletions}: Số dòng code bị xóa đi.
    \item \textbf{diff\_content}: Nội dung chi tiết các đoạn code thay đổi (patch/hunk text) dùng làm input cho AI.
    \item \textbf{created\_at}: Thời điểm lưu bản ghi diff vào hệ thống.
\end{itemize}


% =========================================================
% FILE
% =========================================================

\subsection{File}

\textbf{File (Thực thể yếu -- phụ thuộc Repository):} Cấu trúc tập tin mã nguồn trong kho chứa.

\subsubsection*{Thuộc tính Đơn (Simple / Atomic Attribute)}

\begin{itemize}
    \item \textbf{file\_id (PK)}: Định danh nội bộ của tập tin trong hệ thống.
    \item \textbf{file\_path}: Đường dẫn đầy đủ của file tính từ gốc repository (VD: src/controllers/user\_controller.ts).
    \item \textbf{file\_type}: Loại hoặc đuôi mở rộng của tệp (VD: ts, java, py, md).
    \item \textbf{status}: Trạng thái của file trong cây thư mục (ACTIVE, DELETED, RENAMED).
    \item \textbf{created\_at}: Thời điểm file lần đầu được ghi nhận vào hệ thống.
    \item \textbf{updated\_at}: Thời điểm thông tin file hoặc nội dung file thay đổi gần nhất.
\end{itemize}


% =========================================================
% CODE ENTITY
% =========================================================

\subsection{CodeEntity}

\textbf{CodeEntity (Thực thể yếu -- phụ thuộc File):} Đơn vị mã nguồn cụ thể thu được từ AST (Class, Method, Function,...).

\subsubsection*{Thuộc tính Phức hợp (Composite Attribute)}

\begin{itemize}
    \item \textbf{position}: Vị trí của thực thể trong tập tin, gồm:
    \begin{itemize}
        \item \textbf{start\_line}: Dòng bắt đầu của thực thể code.
        \item \textbf{end\_line}: Dòng kết thúc của thực thể code.
    \end{itemize}
\end{itemize}

\subsubsection*{Thuộc tính Đơn (Simple / Atomic Attribute)}

\begin{itemize}
    \item \textbf{entity\_id (PK)}: Định danh nội bộ của thực thể code.
    \item \textbf{entity\_name}: Tên của thành phần code (VD: processPayment, UserController).
    \item \textbf{entity\_type}: Loại thành phần code (CLASS, METHOD, FUNCTION, INTERFACE, ENDPOINT).
    \item \textbf{signature}: Cấu trúc chữ ký (VD: public void processPayment(String orderId, Double amount)).
    \item \textbf{doc\_comment}: Đoạn comment/docstring nguyên bản đính kèm hàm/lớp (nếu có).
    \item \textbf{created\_at}: Thời điểm thực thể code được bóc tách và lưu vào hệ thống.
    \item \textbf{updated\_at}: Thời điểm cấu trúc của thực thể code thay đổi gần nhất.
\end{itemize}


% =========================================================
% DEPENDENCY
% =========================================================

\subsection{Dependency}

\textbf{Dependency (Thực thể mối kết hợp):} Mối quan hệ phụ thuộc giữa các CodeEntity.

\subsubsection*{Thuộc tính Đơn (Simple / Atomic Attribute)}

\begin{itemize}
    \item \textbf{dependency\_id (PK)}: Định danh duy nhất cho bản ghi phụ thuộc.
    \item \textbf{dependency\_type}: Loại quan hệ phụ thuộc (CALLS, INHERITS, IMPLEMENTS, USES, IMPORTS).
    \item \textbf{created\_at}: Thời điểm hệ thống phân tích AST và ghi nhận mối quan hệ phụ thuộc.
\end{itemize}


% =========================================================
% C. NHÓM KHUÔN MẪU TÀI LIỆU
% =========================================================

\section{Nhóm Khuôn mẫu Tài liệu (Template Engine Domain)}

\subsection{Template}

\textbf{Template (Thực thể mạnh):} Khuôn mẫu chuẩn để AI/User sinh tài liệu.

\subsubsection*{Thuộc tính Đơn (Simple / Atomic Attribute)}

\begin{itemize}
    \item \textbf{template\_id (PK)}: Định danh duy nhất cho khuôn mẫu tài liệu.
    \item \textbf{name}: Tên khuôn mẫu (VD: API Specification Template, System Architecture Doc).
    \item \textbf{description}: Mô tả mục đích sử dụng và phạm vi áp dụng của khuôn mẫu.
    \item \textbf{content\_structure}: Nội dung cấu trúc dàn ý mẫu/Prompt Markdown hướng dẫn AI và User.
    \item \textbf{creator\_type}: Đối tượng khởi tạo khuôn mẫu (SYSTEM, USER, AI).
    \item \textbf{status}: Trạng thái sử dụng (DRAFT, ACTIVE, ARCHIVED).
    \item \textbf{created\_at}: Thời điểm tạo khuôn mẫu.
    \item \textbf{updated\_at}: Thời điểm cập nhật dàn ý/nội dung khuôn mẫu gần nhất.
\end{itemize}


% =========================================================
% TEMPLATE SECTION
% =========================================================

\subsection{TemplateSection}

\textbf{TemplateSection (Thực thể yếu -- phụ thuộc Template):} Mục con/heading trong cấu trúc dàn ý của Template.

\subsubsection*{Thuộc tính Đơn (Simple / Atomic Attribute)}

\begin{itemize}
    \item \textbf{section\_id (PK)}: Định danh duy nhất cho mục dàn ý.
    \item \textbf{heading\_title}: Tiêu đề của mục (VD: 1. API Overview, 2. Request Payload).
    \item \textbf{order\_index}: Thứ tự sắp xếp của mục trong Template.
    \item \textbf{prompt\_instruction}: Câu lệnh hướng dẫn riêng cho AI khi sinh nội dung cho mục này.
    \item \textbf{created\_at}: Thời điểm khởi tạo mục dàn ý.
    \item \textbf{updated\_at}: Thời điểm cập nhật tiêu đề hoặc chỉ thị gần nhất.
\end{itemize}


% =========================================================
% PLACEHOLDER
% =========================================================

\subsection{Placeholder}

\textbf{Placeholder (Thực thể yếu -- phụ thuộc TemplateSection):} Biến/điểm chèn trong Section, làm nhiệm vụ nhận dữ liệu từ code.

\subsubsection*{Thuộc tính Đơn (Simple / Atomic Attribute)}

\begin{itemize}
    \item \textbf{placeholder\_id (PK)}: Định danh duy nhất cho điểm chèn biến.
    \item \textbf{variable\_name}: Mã nhận diện biến (VD: \{\{api\_endpoint\}\}, \{\{request\_schema\}\}).
    \item \textbf{expected\_data\_type}: Loại dữ liệu kỳ vọng từ code (STRING, JSON\_SCHEMA, MERMAID\_DIAGRAM, CODE\_SNIPPET).
    \item \textbf{description}: Mô tả ý nghĩa và ngữ cảnh dữ liệu cần điền vào biến.
    \item \textbf{is\_required}: Cờ xác định biến này bắt buộc AI phải lấp đầy hay không (true/false).
    \item \textbf{created\_at}: Thời điểm tạo điểm chèn biến.
    \item \textbf{updated\_at}: Thời điểm cập nhật thông tin biến gần nhất.
\end{itemize}


% =========================================================
% TEMPLATE MAPPING
% =========================================================

\subsection{TemplateMapping}

\textbf{TemplateMapping (Thực thể mối kết hợp):} Ánh xạ trực tiếp giữa Placeholder và CodeEntity để AI trích xuất thông tin.

\subsubsection*{Thuộc tính Đơn (Simple / Atomic Attribute)}

\begin{itemize}
    \item \textbf{mapping\_id (PK)}: Định danh duy nhất cho bản ghi ánh xạ trích xuất.
    \item \textbf{extraction\_rule}: Quy tắc/Prompt chỉ thị chi tiết để AI biết cách lọc, biến đổi và trích xuất dữ liệu từ CodeEntity (VD: Extract Javadoc + Method Parameters as JSON).
    \item \textbf{is\_active}: Trạng thái hiệu lực của quy tắc ánh xạ (true/false).
    \item \textbf{created\_at}: Thời điểm khởi tạo quy tắc ánh xạ.
    \item \textbf{updated\_at}: Thời điểm cập nhật quy tắc trích xuất gần nhất.
\end{itemize}


% =========================================================
% D. NHÓM TÀI LIỆU & VÒNG ĐỜI DỮ LIỆU
% =========================================================

\section{Nhóm Tài liệu \& Vòng đời Dữ liệu (Documentation Domain)}

\subsection{Document}

\textbf{Document (Thực thể mạnh):} Tài liệu gốc (README, API Reference, Architecture Guide, ADR,...).

\subsubsection*{Thuộc tính Đơn (Simple / Atomic Attribute)}

\begin{itemize}
    \item \textbf{document\_id (PK)}: Định danh duy nhất cho tài liệu chính thức.
    \item \textbf{title}: Tên/Tiêu đề của tài liệu (VD: API Reference v2, System Architecture).
    \item \textbf{document\_type}: Phân loại loại tài liệu (README, API\_REFERENCE, ARCHITECTURE\_GUIDE, ADR, CUSTOM).
    \item \textbf{status}: Trạng thái vòng đời của tài liệu (PUBLISHED, ARCHIVED, DEPRECATED).
    \item \textbf{created\_at}: Thời điểm khởi tạo tài liệu.
    \item \textbf{updated\_at}: Thời điểm cập nhật tài liệu gần nhất.
    \item \textbf{deleted\_at}: Thời điểm xóa mềm tài liệu (Nullable).
\end{itemize}


% =========================================================
% DOCUMENT VERSION
% =========================================================

\subsection{DocumentVersion}

\textbf{DocumentVersion (Thực thể yếu -- phụ thuộc Documentation):} Bản ghi nội dung tài liệu theo từng mốc thời gian và trạng thái (Draft, In Review, Approved, Published).

\subsubsection*{Thuộc tính Đơn (Simple / Atomic Attribute)}

\begin{itemize}
    \item \textbf{version\_id (PK)}: Định danh duy nhất cho bản ghi phiên bản.
    \item \textbf{version\_number}: Số hiệu phiên bản tài liệu (VD: 1.0.0, 1.1.0, 2.0-draft).
    \item \textbf{content}: Toàn bộ nội dung văn bản tài liệu dạng Markdown/HTML tại thời điểm chụp phiên bản.
    \item \textbf{change\_summary}: Tóm tắt mô tả lý do hoặc nội dung những điểm thay đổi trong phiên bản này.
    \item \textbf{status}: Trạng thái vòng đời phiên bản (DRAFT, IN\_REVIEW, APPROVED, PUBLISHED).
    \item \textbf{created\_at}: Thời điểm khởi tạo bản ghi phiên bản.
    \item \textbf{updated\_at}: Thời điểm cập nhật trạng thái phiên bản gần nhất.
\end{itemize}


% =========================================================
% CHANGE LOG
% =========================================================

\subsection{ChangeLog}

\textbf{ChangeLog (Thực thể yếu -- phụ thuộc DocumentVersion):} Lịch sử vết thay đổi chi tiết giữa các phiên bản tài liệu.

\subsubsection*{Thuộc tính Đơn (Simple / Atomic Attribute)}

\begin{itemize}
    \item \textbf{log\_id (PK)}: Định danh duy nhất cho bản ghi nhật ký thay đổi.
    \item \textbf{change\_type}: Loại thao tác biến đổi (ADDED, MODIFIED, DELETED, RESTORED).
    \item \textbf{change\_description}: Mô tả tóm tắt ý nghĩa sự thay đổi.
    \item \textbf{diff\_content}: Nội dung chênh lệch chi tiết (dạng Git patch / JSON delta) để dựng giao diện so sánh code/doc.
    \item \textbf{created\_at}: Thời điểm ghi nhận vết thay đổi.
    \item \textbf{updated\_at}: Thời điểm cập nhật bản ghi nhật ký (nếu có bổ sung metadata).
\end{itemize}


% =========================================================
% DOCUMENT ASSET
% =========================================================

\subsection{DocumentAsset}

\textbf{DocumentAsset (Thực thể yếu -- phụ thuộc DocumentVersion):} Tài nguyên đa phương tiện (hình ảnh, sơ đồ) gắn kèm tài liệu.

\subsubsection*{Thuộc tính Đơn (Simple / Atomic Attribute)}

\begin{itemize}
    \item \textbf{asset\_id (PK)}: Định danh duy nhất cho tài nguyên đa phương tiện.
    \item \textbf{file\_name}: Tên file gốc hoặc tên sơ đồ (VD: architecture\_v1.png, sequence\_payment.svg).
    \item \textbf{asset\_url}: Đường dẫn công khai/CDN đến tệp lưu trữ đa phương tiện.
    \item \textbf{asset\_type}: Phân loại định dạng tài nguyên (IMAGE, DIAGRAM, ARCH\_SCHEMA, ATTACHMENT).
    \item \textbf{created\_at}: Thời điểm tải lên hoặc xuất sơ đồ vào hệ thống.
    \item \textbf{updated\_at}: Thời điểm cập nhật thông tin tài nguyên gần nhất.
\end{itemize}


% =========================================================
% TICKET REFERENCE
% =========================================================

\subsection{TicketReference}

\textbf{TicketReference (Thực thể phụ thuộc):} Mã tham chiếu liên kết đến các công cụ quản lý công việc bên ngoài (Jira Ticket ID, GitHub Issue).

\subsubsection*{Thuộc tính Đơn (Simple / Atomic Attribute)}

\begin{itemize}
    \item \textbf{reference\_id (PK)}: Định danh duy nhất cho bản ghi tham chiếu ticket.
    \item \textbf{platform}: Tên hệ thống quản lý công việc bên ngoài (JIRA, GITHUB\_ISSUE, LINEAR, TRELLO).
    \item \textbf{ticket\_code}: Mã định danh ticket trên hệ thống gốc (VD: PROJ-123, \#45).
    \item \textbf{created\_at}: Thời điểm tạo bản ghi liên kết.
    \item \textbf{updated\_at}: Thời điểm cập nhật liên kết gần nhất.
\end{itemize}


% =========================================================
% E. NHÓM KIỂM DUYỆT & PHÊ DUYỆT
% =========================================================

\section{Nhóm Kiểm duyệt \& Phê duyệt (Review \& Governance Domain)}

\subsection{Review}

\textbf{Review (Thực thể mối kết hợp/Sự kiện):} Tiến trình Staff thực hiện đánh giá kỹ thuật trên một DocumentVersion.

\subsubsection*{Thuộc tính Đơn}

\begin{itemize}
    \item \textbf{review\_id (PK)}: Định danh duy nhất cho phiên đánh giá kỹ thuật.
    \item \textbf{status}: Trạng thái của tiến trình review (PENDING, IN\_PROGRESS, COMPLETED).
    \item \textbf{decision}: Kết quả thẩm định của reviewer (APPROVED, REJECTED, CHANGES\_REQUESTED).
    \item \textbf{summary\_notes}: Nhận xét hoặc ghi chú tổng quan của reviewer (Nullable).
    \item \textbf{created\_at}: Thời điểm bắt đầu phiên đánh giá.
    \item \textbf{updated\_at}: Thời điểm chốt kết quả hoặc cập nhật tiến trình review.
\end{itemize}


% =========================================================
% REVIEW COMMENT
% =========================================================

\subsection{ReviewComment}

\textbf{ReviewComment (Thực thể yếu -- phụ thuộc Review):} Phản hồi, góp ý chi tiết trong quá trình kiểm duyệt.

\subsubsection*{Thuộc tính Đơn}

\begin{itemize}
    \item \textbf{comment\_id (PK)}: Định danh duy nhất cho bản ghi bình luận/góp ý.
    \item \textbf{content}: Nội dung chi tiết của lời nhận xét, phản hồi.
    \item \textbf{is\_resolved}: Trạng thái xử lý góp ý (false: Chưa sửa / true: Đã sửa xong).
    \item \textbf{created\_at}: Thời điểm gửi bình luận.
    \item \textbf{updated\_at}: Thời điểm cập nhật nội dung bình luận gần nhất.
\end{itemize}


% =========================================================
% APPROVAL
% =========================================================

\subsection{Approval}

\textbf{Approval (Thực thể mối kết hợp/Sự kiện):} Quyết định phê duyệt xuất bản/commit tài liệu chính thức từ Manager.

\subsubsection*{Thuộc tính Đơn}

\begin{itemize}
    \item \textbf{approval\_id (PK)}: Định danh duy nhất cho bản ghi quyết định phê duyệt.
    \item \textbf{decision}: Trạng thái/Kết quả quyết định (APPROVED, REJECTED).
    \item \textbf{rejection\_reason}: Lý do từ chối xuất bản (Nullable, dùng khi decision = REJECTED).
    \item \textbf{approved\_at}: Mốc thời gian ra quyết định phê duyệt hoặc từ chối.
\end{itemize}


% =========================================================
% F. NHÓM AI EVIDENCE, GIÁM SÁT & KIỂM TOÁN
% =========================================================

\section{Nhóm AI Evidence, Giám sát \& Kiểm toán (AI Evidence \& Observability Domain)}

\subsection{GroundingEvidence}

\textbf{GroundingEvidence (Thực thể yếu -- phụ thuộc DocumentVersion):} Bảng minh chứng chứa liên kết nguồn gốc (File, CodeEntity, Commit) AI dùng để tạo nội dung, phục vụ chống hallucination.

\subsubsection*{Thuộc tính Đơn}

\begin{itemize}
    \item \textbf{evidence\_id (PK)}: Định danh duy nhất cho bản ghi bằng chứng minh chứng.
    \item \textbf{snippet\_content}: Đoạn mã nguồn hoặc văn bản thô trích xuất từ code làm căn cứ trực tiếp cho AI sinh nội dung.
    \item \textbf{purpose}: Phân loại mục đích lưu vết (HALLUCINATION\_PREVENTION, CODE\_REFERENCE, AST\_VERIFICATION).
    \item \textbf{created\_at}: Thời điểm ghi nhận bản ghi bằng chứng.
\end{itemize}


% =========================================================
% DRIFT ALERT
% =========================================================

\subsection{DriftAlert}

\textbf{DriftAlert (Thực thể sự kiện):} Cảnh báo phát hiện tài liệu bị lệch ngữ nghĩa/lỗi thời so với code mới.

\subsubsection*{Thuộc tính Khóa (Key Attribute)}

\begin{itemize}
    \item \textbf{alert\_id}: Định danh duy nhất cho từng bản ghi cảnh báo.
\end{itemize}

\subsubsection*{Thuộc tính Đơn / Nguyên tố (Simple / Atomic Attributes)}

\begin{itemize}
    \item \textbf{drift\_type}: Phân loại hình thức lệch pha (SEMANTIC, OUTDATED).
    \item \textbf{status}: Trạng thái xử lý cảnh báo (OPEN, RESOLVED, IGNORED).
    \item \textbf{detected\_at}: Thời điểm hệ thống AI quét và phát hiện ra độ lệch.
\end{itemize}


% =========================================================
% REPORT
% =========================================================

\subsection{Report}

\textbf{Report (Thực thể nghiệp vụ):} Dữ liệu báo cáo tổng quan về chất lượng tài liệu, tình trạng drift và chỉ số hệ thống.

\subsubsection*{Thuộc tính Đơn}

\begin{itemize}
    \item \textbf{report\_id (PK)}: Định danh duy nhất cho bản ghi báo cáo.
    \item \textbf{report\_type}: Phân loại báo cáo (DAILY, WEEKLY, MONTHLY).
    \item \textbf{total\_documents}: Tổng số tài liệu hiện có trong workspace.
    \item \textbf{drifted\_documents\_count}: Số lượng tài liệu đang bị cảnh báo lệch pha/lỗi thời.
    \item \textbf{ai\_generations\_count}: Tổng số lượt gọi AI sinh tài liệu trong kỳ.
    \item \textbf{generated\_at}: Mốc thời gian chốt số liệu báo cáo.
\end{itemize}


% =========================================================
% AUDIT LOG
% =========================================================

\subsection{AuditLog}

\textbf{AuditLog (Thực thể nhật ký):} Lưu vết kiểm toán toàn bộ hành vi quản trị và thay đổi quan trọng trong hệ thống.

\subsubsection*{Thuộc tính Đơn}

\begin{itemize}
    \item \textbf{log\_id (PK)}: Định danh duy nhất cho bản ghi nhật ký kiểm toán.
    \item \textbf{action}: Tên thao tác thực hiện (UPDATE\_ROLE, DELETE\_DOC, CHANGE\_POLICY, DISABLE\_USER).
    \item \textbf{target\_entity}: Tên bảng/loại đối tượng bị tác động (Role, Workspace, Document).
    \item \textbf{target\_entity\_id}: Khóa chính của đối tượng bị tác động.
    \item \textbf{details}: Chi tiết biến động dữ liệu ngắn gọn (Nullable).
    \item \textbf{created\_at}: Thời điểm xảy ra sự kiện kiểm toán.
\end{itemize}


% =========================================================
% AI JOB
% =========================================================

\subsection{AIJob}

\textbf{AIJob (Thực thể tác vụ/tiến trình):} Tiến trình xử lý ngầm của AI (generate, update, detect drift, indexing, RAG).

\subsubsection*{1. Thuộc tính Khóa (Key Attribute)}

\begin{itemize}
    \item \textbf{job\_id (PK)}: Định danh duy nhất cho từng tiến trình xử lý ngầm.
\end{itemize}

\subsubsection*{2. Thuộc tính Đơn (Atomic Attributes)}

\begin{itemize}
    \item \textbf{job\_type}: Phân loại tác vụ AI (GENERATE\_DOC, UPDATE\_AST, DETECT\_DRIFT, INDEX\_RAG, SYNC\_REPO).
    \item \textbf{status}: Trạng thái vòng đời của job (PENDING, PROCESSING, COMPLETED, FAILED).
    \item \textbf{target\_entity\_type}: Tên loại đối tượng mục tiêu (Document, PullRequest, Repository, DocumentDraft).
    \item \textbf{target\_entity\_id}: Khóa chính của đối tượng mục tiêu đang được xử lý.
    \item \textbf{error\_message}: Dòng thông báo lỗi chi tiết khi trạng thái là FAILED (Nullable).
\end{itemize}

\subsubsection*{3. Thuộc tính Thời gian (Temporal Attributes)}

\begin{itemize}
    \item \textbf{created\_at}: Thời điểm job được đưa vào hàng đợi (Queue).
    \item \textbf{started\_at}: Thời điểm worker bắt đầu thực thi job.
    \item \textbf{finished\_at}: Thời điểm job hoàn thành hoặc bị ngắt do lỗi (Nullable).
\end{itemize}

\subsubsection*{4. Thuộc tính Dẫn xuất (Derived Attribute)}

\begin{itemize}
    \item \textbf{execution\_time}: Thời gian xử lý thực tế của job, không lưu cố định trong CSDL, được tính bằng:
    \[
        execution\_time = finished\_at - started\_at
    \]
\end{itemize}